\documentclass[11pt]{article}
\usepackage{amsmath,amssymb,amsthm,booktabs,geometry,url}
\geometry{margin=1in}
\setlength{\emergencystretch}{1em}
\newtheorem{theorem}{Theorem}
\newtheorem{lemma}[theorem]{Lemma}
\newtheorem{corollary}[theorem]{Corollary}
\newcommand{\im}{\operatorname{im}}
\newcommand{\rank}{\operatorname{rank}}
\newcommand{\wt}{\operatorname{wt}}
\title{Intersection counts and orthogonal graph codes}
\author{Working manuscript: distinct review pending}
\date{}
\begin{document}\maketitle

\paragraph{Attribution and scope.}
Counting common neighbours in an association scheme is classical. For
uniform subset graphs the count below is Proposition~3 of Key and Rodrigues,
\emph{LCD codes from adjacency matrices of graphs}, manuscript dated
4 August 2017 (F017). Their examples and classifications are recovered here
from one scalar test and the common orthogonal projection theorem. We also
give uniform proofs for the Hamming examples in their conclusion, correcting
the stated complement-dual identity for even binary cubes. No historical
novelty or resolution of an open conjecture is claimed. Dimensions and
distances quoted from other papers in F017 require their own proofs or
explicitly retained inputs; the scalar test alone does not prove them.

Write $C_M=\im M$ for a symmetric matrix $M$ over $\mathbb F_p$. Let $A$
be the adjacency matrix of a regular simple graph on $N$ vertices, of degree
$\nu$, and put $B=J-I-A$. Following F017, call $C_A$ RLCD when
$C_A=C_{A+I}^{\perp}$, and CLCD when $C_A=C_B^{\perp}$.
Suppose ordered pairs of vertices are partitioned into classes, indexed
by $j$, on which both the adjacency indicator $a_j$ and the number $i_j$
of common neighbours are constant. Include the diagonal class among them.

\begin{theorem}\label{thm:test}
Put $P_j=i_j+a_j$ and $\bar\nu=N-1-\nu$. Then:
\begin{enumerate}
\item $C_A$ is RLCD if and only if $P_j=0$ in $\mathbb F_p$ for every $j$.
\item $C_A\perp C_B$ if and only if $\nu-P_j=0$ for every $j$.
In this event both codes are LCD. They are complementary if and only if
$(\nu,\bar\nu)\ne(0,0)$ in the field; otherwise
$\mathbb F_p^N=C_A\perp\!\oplus C_B\perp\!\oplus\langle\mathbf1\rangle$.
\item $C_B$ is RLCD if and only if $\bar\nu-\nu+P_j=0$ for every $j$.
\end{enumerate}
\end{theorem}
\begin{proof}
The entries of $A^2+A$ are $P_j$. Thus their vanishing is exactly
$A(A+I)=0$. The matrices $-A$ and $I+A$ are then complementary symmetric
idempotents. Conversely the asserted duality implies that all their row
inner products vanish. Next
$AB=AJ-A-A^2=\nu J-A-A^2$, proving the orthogonality test.
The remaining claims in part~2 are the common projection theorem, applied
to this actual pair of matrices; no condition on their complex spectra is
needed. Finally $B(B+I)=B(J-A)=\bar\nu J-AB$, which gives part~3 by the
same reflexive argument.
\end{proof}

\paragraph{Uniform subset graphs.}
Let $n\geq2k+1$, $k\geq2$ and $0\leq r<k$. The vertices of $\Gamma(n,k,r)$
are the $k$-subsets of an $n$-element set; adjacency means intersection of
size $r$. For two vertices meeting in $j$ elements, the common-neighbour
count is
\begin{equation}\label{eq:subset}
 i_j=\sum_a\binom ja\binom{k-j}{r-a}^{\!2}
                  \binom{n-2k+j}{k-2r+a},\qquad 0\leq j\leq k.
\end{equation}
Here binomial coefficients with invalid lower argument are zero.
To prove the formula, partition the ground set into the intersection,
the two differences and the outside. A common neighbour chooses $a$
elements from the intersection, $r-a$ from each difference and the remaining
$k-2r+a$ outside. These choices are disjoint and exhaustive.
In particular $i_k=\nu=\binom kr\binom{n-k}{k-r}$ and
$P_j=i_j+\mathbf1_{j=r}$, so Theorem~\ref{thm:test} applies directly.

For the Johnson graph $r=k-1$, the only nonzero counts off the diagonal
are $i_{k-2}=4$ and $i_{k-1}=n-2$. Thus RLCD holds exactly for $p=2$ and
odd $n$. Under those conditions $\nu=0$, so CLCD holds exactly when
$\binom nk$ is even. For $k\geq3$, complement orthogonality itself forces
$\nu=0$ from $P_0=0$;
then $P_{k-2}=4$ and $P_{k-1}=n-1$ force the same characteristic and parity.
For $k=2$, subtraction of the two off-diagonal tests yields $n=5$ in the
field, while $\nu=4$ gives $2(n-2)=4$, hence characteristic two and odd $n$.
The binary dimensions and minimum words follow from the separate simplex
boundary theorem, not from this count alone.

The domain qualification matters. F017 Corollary~1(1) does not state it
explicitly. For example $J(4,3)=K_4$ is RLCD in characteristic three.
More generally the adjacency matrix of $J(n,n-1)=K_n$, for $n\geq3$, is
$J-I$, so $(J-I)J=(n-1)J$: it is RLCD exactly when $p\mid(n-1)$.
Its complement is empty, so it is CLCD exactly when $J-I$ is invertible,
equivalently $p\nmid(n-1)$. Thus the unrestricted source exclusion is false.
Formula~\eqref{eq:subset} itself remains valid for $n\geq k$ when only
realized intersections $\max(0,2k-n)\leq j\leq k$ are tested; the table
below retains its stated $n\geq2k+1$ domain.

For a Kneser graph ($r=0$), formula~\eqref{eq:subset} reduces to
$i_j=\binom{n-2k+j}{k}$. If $n=2k+1$, then $i_0=0$ and $P_0=1$, so the
code is never RLCD. For $k\geq3$ it is not CLCD either: complement
orthogonality would force both $\nu=P_0=1$ and $\nu=i_1=0$.
For $k=2$ there is an exception. The Petersen graph has $(N,\nu)=(10,3)$,
$P_0=1$, $P_1=1$ and $P_2=3$. Complement orthogonality holds exactly
in characteristic two, where $\nu=1$ and $\bar\nu=0$; hence this code
is CLCD. F017 Corollary~1(2) states the exclusion for every $k$.
Its CLCD assertion requires this endpoint correction. Its RLCD exclusion
is valid for every $k\geq2$.

\begin{table}[ht]
\centering\small
\begin{tabular}{ccp{0.32\textwidth}p{0.32\textwidth}}
\toprule
$k$&$r$&RLCD&CLCD\\\midrule
2&0&$p=2$, $n\equiv3\pmod4$&$p=2$, $n\equiv1\pmod4$\\
3&0&$p=2$, $n\equiv1\pmod4$; $p=3$, $n\equiv5\pmod9$;
$p\geq5$, $n\equiv5\pmod p$&$p=2$, $n\equiv1\pmod4$;
$p=3$, $n\equiv2,8\pmod9$; $p\geq5$, $n\equiv5\pmod p$\\
4&0&$p=2$, $n\equiv7\pmod8$&$p=2$, $n\equiv3\pmod8$\\
3&1&$p=2$, $n\equiv0\pmod4$&$p=2$, $n\equiv0,1\pmod4$\\
4&1&$p=2$, $n\equiv1\pmod4$; $p=3$, $n\equiv6\pmod9$
&$p=2$, $n\equiv1\pmod8$; $p=3$, $n\equiv6\pmod9$\\
\bottomrule
\end{tabular}
\caption{Complete parameter conditions for the five families specified in
F017, with $n\geq2k+1$. Each entry is an if and only if statement.}
\label{tab:families}
\end{table}

\paragraph{Completeness of the arithmetic classification.}
Here is a finite exact certificate for Table~\ref{tab:families}, including
the exclusion of all other primes. For each fixed $(k,r)$, regard the
$P_j$ and $\nu-P_j$ as polynomials in $n$ over $\mathbb Q$ using
\eqref{eq:subset}. For each list $f_0,\ldots,f_k$, the accompanying program
records polynomial coefficients $b_j$ and a monic polynomial $g$ such that
\[
 \sum_j b_jf_j=g,\qquad g\mid f_j\quad\text{for all }j.
\]
Both assertions are verified by exact polynomial arithmetic. Let $D$ be
the product of a common denominator of all coefficients of all $b_j$ and
a common denominator of all coefficients of all $f_j$.
For $p\nmid D$, the simultaneous congruences are equivalent to $g(n)=0$.
In the two $k=3,r=0$ lists, $g=n-5$; in all other lists, $g=1$.
The resulting exceptional primes, including primes subsequently excluded
by the finite tests, are
\[
\begin{array}{c|cc}
(k,r)&P_j&\nu-P_j\\\hline
(2,0)&2&2\\
(3,0)&2,3&2,3\\
(4,0)&2,3&2,3,17\\
(3,1)&2,3,5&2,3,23\\
(4,1)&2,3,5&2,3,5.
\end{array}
\]
For an exceptional prime choose the least power $p^a>k$. For $0\leq b\leq k$,
the identity $(1+X)^{p^a}=1+X^{p^a}$ over $\mathbb F_p$ proves
$\binom{m+p^a}{b}=\binom mb$ modulo $p$ for nonnegative $m$.
Consequently all the counts, the degree and $\binom nk$ have period $p^a$
modulo $p$ in the domain $n\geq2k+1$. Testing one admissible lift of every
residue modulo $p^a$ is therefore exhaustive, not a bounded experiment.
The program records these residues and applies the constant-line exclusion
in Theorem~\ref{thm:test} to obtain CLCD. For $k=3,r=0,p\geq5,n\equiv5$,
the degree is zero but $\binom n3-1\equiv9\ne0$, so no further exception
occurs. This proves the table from its exact arithmetic certificate.

\paragraph{Hamming graphs.}
Let $H_k(d,q)$ have vertex set $\{0,\ldots,q-1\}^d$, with adjacency at
Hamming distance $k$, where $1\leq k\leq d$ and $q\geq2$. If two vertices
have distance $j$, their common-neighbour count is
\begin{equation}\label{eq:hamming}
 i_j=\sum_a\binom{d-j}{a}(q-1)^a
 \frac{j!}{b!^2e!}(q-2)^e,
 \qquad b=j+a-k,\quad e=2k-j-2a,
\end{equation}
where terms with negative factorial arguments are omitted. In their $d-j$
equal coordinates, choose the $a$ positions where a common neighbour differs.
Among the $j$ different coordinates, let $b$ agree with the first vertex,
$b$ with the second, and $e$ with neither. The distance equations give the
stated values of $b,e$, with $2b+e=j$. Multinomial counting proves
\eqref{eq:hamming}. Theorem~\ref{thm:test} now applies with
$N=q^d$, $\nu=i_0=\binom dk(q-1)^k$ and $P_j=i_j+\mathbf1_{j=k}$.

\begin{corollary}\label{cor:oddq}
For odd $q$, every binary distance-$k$ Hamming code is RLCD.
For the distance-two binary Hamming graph over $\mathbb F_3$, RLCD holds
whenever $d\equiv0\pmod3$.
\end{corollary}
\begin{proof}
Let $X_i$ act as the complete-graph adjacency matrix $J_q-I_q$ in coordinate
$i$ and as the identity in all others. Over $\mathbb F_2$, odd $q$ gives
$X_i^2=X_i$, and these matrices commute. The distance-$k$ adjacency matrix
is $A_k=\sum_{|S|=k}\prod_{i\in S}X_i$. Squaring cancels mixed terms in
characteristic two, so $A_k^2=A_k$ and $A_k(A_k+I)=0$.
For $q=2,k=2$, the only possible common-neighbour distances are $0,2,4$,
with counts $\binom d2$, $2(d-2)$ and $6$, respectively. If $3\mid d$,
these are $0,-1,0$ modulo three. This is precisely the RLCD test.
\end{proof}

For completeness, the Hamming spectrum used to compute the dimensions
comes from the same tensor construction. Over $\mathbb Q$, $J_q-I_q$
has eigenvalue $q-1$ on the constant vectors and $-1$ on their complement.
Choosing $t$ nonconstant tensor factors therefore gives
\[
 \lambda_t=\sum_i(-1)^i(q-1)^{k-i}\binom ti\binom{d-t}{k-i},
 \qquad m_t=\binom dt(q-1)^t.
\]
The common LCD rank lemma sums precisely those multiplicities for which
$p\nmid\lambda_t$. In particular the ternary $H_2(6,2)$ and $H_2(9,2)$
codes have dimensions $42$ and $342$. This rank step is justified by the
proved LCD property, rather than an assumption of diagonalizability after
reduction modulo $p$.

\begin{theorem}\label{thm:cube}
Over $\mathbb F_2$, the distance-$k$ binary Hamming adjacency matrix satisfies
$A_k^2=\binom dk I$. For positive even $d$, the code of the ordinary cube
$Q_d=H_1(d,2)$ is self-dual of dimension $2^{d-1}$. Its complement adjacency
matrix is invertible, so the code is not the dual of the complement code.
\end{theorem}
\begin{proof}
Index coordinates by $\mathbb F_2^d$. Write $T_s$ for translation by $s$.
Then $A_k=\sum_{\wt(s)=k}T_s$. The translations commute, each satisfies
$T_s^2=I$, and cross terms cancel in the square, proving the first assertion.
Write $M_t$ for the adjacency matrix of the ordinary cube $Q_t$.
Order cube vertices by their last coordinate. If $d$ is even,
\[
 M_d=\begin{pmatrix}M_{d-1}&I\\I&M_{d-1}\end{pmatrix},\qquad
 M_{d-1}^2=I.
\]
The kernel consists exactly of $(x,M_{d-1}x)$, so both nullity and rank are
$2^{d-1}$. Since $M_d^2=0$ and $M_d$ is symmetric, its image equals its
orthogonal complement. On the other hand, with $B=J+I+M_d$ we have
$J^2=2^dJ=0$ and $JM_d=M_dJ=dJ=0$, so $B^2=I$.
\end{proof}

The conclusion of F017 displays the complement-dual identity for
$Q_6,Q_8,Q_{10}$. Direct reading of the formula confirms the complement bar;
Theorem~\ref{thm:cube} disproves that identity and supplies the corrected
self-duality statement. The two remaining Hamming examples in that
conclusion are instances of Corollary~\ref{cor:oddq} and are retained.

\paragraph{Verification and limits.}
The required scalar certificate is \path{check_association_counts.py} and
its exact polynomial and residue report. Separate explicit finite graphs
test the counting formulas on all vertex pairs in those test graphs;
these tests support implementation correctness and are not used to infer
the infinite classification. No finite minimum-distance assertion is
inferred from this manuscript. The characteristic-polynomial rank lemma
and the general projection theorem are shared dependencies, each charged
in the common theory when this work is measured.
\end{document}
